\section{Week 2: Optimization, Generalization, and Scaling Intuition}

\subsection{Learning Objectives}

By the end of this week, you should be able to:
\begin{itemize}
  \item Explain why gradient-based optimization behaves counterintuitively in deep networks
  \item Distinguish optimization failure from generalization failure
  \item Understand why overparameterization often \emph{helps} generalization
  \item Reason about learning rates, optimizers, and schedules as system-level design choices
  \item Connect small-scale experiments to large-scale LLM training behavior
\end{itemize}

\subsection{LLM Components Covered}

\textbf{Training dynamics}
\begin{itemize}
  \item Gradient descent and stochasticity
  \item Adam / AdamW vs SGD
  \item Learning-rate schedules
  \item Overparameterization
\end{itemize}

\textbf{Systems relevance}
\begin{itemize}
  \item Training instability and regressions
  \item Sensitivity to hyperparameters
  \item Why scaling laws emerge empirically
\end{itemize}

\subsection{Conceptual Core: Optimization Is Not Inference}

Training an LLM is a large-scale numerical optimization problem, not a logical inference procedure.
The training objective from Week~1 is optimized via variants of stochastic gradient descent (SGD):

\begin{equation}
\theta_{t+1} = \theta_t - \eta_t \nabla_\theta \mathcal{L}(\theta_t)
\end{equation}

where:
\begin{itemize}
  \item $\theta$ are model parameters,
  \item $\eta_t$ is the learning rate,
  \item $\mathcal{L}$ is the empirical cross-entropy loss.
\end{itemize}

Crucially, this optimization occurs in:
\begin{itemize}
  \item very high-dimensional parameter spaces,
  \item with non-convex loss surfaces,
  \item using noisy gradient estimates.
\end{itemize}

Yet, modern deep networks reliably converge.

\subsection{Why Overparameterization Does Not Doom Generalization}

Classical intuition suggests that models with more parameters than data should overfit catastrophically.
Empirically, deep networks exhibit \emph{benign overfitting}:
\begin{itemize}
  \item training loss goes to (near) zero,
  \item test loss does not explode,
  \item larger models often generalize better than smaller ones.
\end{itemize}

This behavior is not fully explained theoretically, but several mechanisms are believed to contribute:
\begin{itemize}
  \item implicit regularization from optimization dynamics,
  \item flat minima bias,
  \item data structure and redundancy,
  \item early stopping.
\end{itemize}

For LLMs, this observation underpins the entire scaling paradigm.

\subsection{Optimizers as Design Choices}

\subsubsection{SGD vs AdamW}

SGD updates parameters using raw gradient estimates.
Adam-style optimizers maintain adaptive per-parameter learning rates.

AdamW decouples weight decay from gradient updates, improving stability.

\begin{center}
\begin{tabular}{|l|l|l|}
\hline
Optimizer & Strengths & Risks \\
\hline
SGD & Simplicity, stability & Slow convergence \\
Adam & Fast initial progress & Overfitting, sharp minima \\
AdamW & Better generalization & Sensitive hyperparameters \\
\hline
\end{tabular}
\end{center}

Large language models almost universally rely on AdamW or closely related variants.

\subsection{Learning Rate Schedules}

The learning rate $\eta_t$ is often the single most important hyperparameter.

Common schedules include:
\begin{itemize}
  \item constant learning rate,
  \item linear warmup + decay,
  \item cosine decay.
\end{itemize}

Warmup phases are especially critical for large models to avoid early divergence.

\subsection{Empirical Scaling and Its Consequences}

Empirical scaling laws show that loss decreases smoothly as a function of:
\begin{itemize}
  \item model size,
  \item dataset size,
  \item compute budget.
\end{itemize}

These relationships hold over many orders of magnitude and motivate:
\begin{itemize}
  \item predictability of returns on compute investment,
  \item industrial-scale training pipelines,
  \item architectural choices that prioritize scalability.
\end{itemize}

\subsection{Systems Engineering Implications}

\textbf{Requirements}
\begin{itemize}
  \item Training must be stable and repeatable
  \item Failures must be detectable early
\end{itemize}

\textbf{Architecture}
\begin{itemize}
  \item Gradient clipping
  \item Learning-rate warmup
  \item Checkpointing and rollback
\end{itemize}

\textbf{Verification \& Validation}
\begin{itemize}
  \item Loss curve monitoring
  \item Seed sensitivity testing
  \item Small-scale proxy experiments
\end{itemize}

\textbf{Operations}
\begin{itemize}
  \item Automated regression detection
  \item Cost–quality tradeoff tracking
\end{itemize}

\subsection{Human Factors and Failure Modes}

A recurring human failure mode is attributing training behavior to ``model intelligence''
rather than optimization dynamics.

Examples:
\begin{itemize}
  \item ``The model learned reasoning'' vs ``the loss surface permits generalizing minima''
  \item Overinterpreting early training instability
\end{itemize}

Understanding optimization mechanics is essential for realistic expectations.

\subsection{Key References}

\begin{itemize}
  \item \citet{goodfellow2016dl} — optimization and generalization in deep learning
  \item \citet{zhang2017understanding} — empirical study of overparameterization
\end{itemize}

\subsection{Recommended University Lectures}

\begin{itemize}
  \item MIT 6.S191 — optimization and deep learning practice  
  \url{https://introtodeeplearning.mit.edu}

  \item Berkeley CS182 — deep learning foundations  
  \url{https://cs182.berkeley.edu}

  \item Stanford CS229 — gradient descent and generalization  
  \url{https://cs229.stanford.edu}
\end{itemize}

\subsection{Practitioner Lab}

See \texttt{labs/week02/week02\_lab\_optimization\_generalization.ipynb}.

\subsection{Week 2 Takeaway}

\begin{quote}
Training behavior in LLMs is governed less by ``intelligence'' than by optimization dynamics operating in very high dimensions.
\end{quote}

